Le second degré est le premier endroit du programme où une seule identité gouverne tout le
chapitre : la forme canonique. Le discriminant, les racines, le signe du trinôme, le sommet
de la parabole et son axe de symétrie en sont des lectures.

La démonstration de la somme et du produit des racines mérite un mot. Plutôt que d'identifier
les coefficients, elle évalue l'identité de factorisation en trois points : en $0$ on lit le
produit, et la différence des évaluations en $1$ et en $-1$ élimine deux coefficients et
laisse la somme. C'est plus court, et c'est exactement ce que fait le script formel.

Le chapitre se termine par la fonction homographique, où apparaît pour la première fois une
valeur interdite — le point où le dénominateur s'annule. C'est le premier objet du programme
dont le domaine n'est pas $\mathbb{R}$ tout entier, et cela se paie en hypothèses dans chaque
énoncé.
